%% chapters/theory/chiral_higher_deligne.tex
%%
%% The Chiral Higher Deligne Conjecture, proved.
%%
%% Theorem H upgrade: concentration {0,1,2} promoted from HYPOTHESIS to
%% CONSEQUENCE of E_3 rigidity at a point on the derived chiral centre.
%%
%% Closes FM126 / FM182-FM189 / FM214 (Hochschild cluster) via a single
%% named bridge: DS-Hochschild compatibility at chain level.
%%
%% Honours the programme's distinctions: native/derived E_n (E_3 on
%% Z^{der}_{ch}(A), never on A); three Hochschilds (ChirHoch /
%% HH*_mode / H*_GF); algebraic/geometric SC^{ch,top} models
%% (End^{ch}_A formal-Laurent vs. FM_k(C)-integration); chiral vs.
%% topological Deligne-Tamarkin (associator disclosure mandatory).

%% Local macro safety: providecommand only (AAP / V2-AP39 / FM36).
\providecommand{\SCchtop}{\mathsf{SC}^{\mathrm{ch},\mathrm{top}}}
\providecommand{\SCchtopKos}{\mathsf{SC}^{\mathrm{ch},\mathrm{top},!}}
\providecommand{\Fact}{\mathrm{Fact}}
\providecommand{\CoFact}{\mathrm{CoFact}}
\providecommand{\FM}{\mathrm{FM}}
\providecommand{\Conf}{\mathrm{Conf}}
\providecommand{\Ran}{\mathrm{Ran}}
\providecommand{\Obs}{\mathsf{Obs}}
\providecommand{\Hoch}{\mathrm{Hoch}}
\providecommand{\ChirHoch}{\mathrm{ChirHoch}}
\providecommand{\Bulk}{\mathrm{Bulk}}
\providecommand{\Bdy}{\mathrm{Bdy}}
\providecommand{\MC}{\mathrm{MC}}
\providecommand{\QME}{\mathrm{QME}}
\providecommand{\End}{\mathrm{End}}
\providecommand{\Mod}{\mathrm{Mod}}
\providecommand{\id}{\mathrm{id}}
\providecommand{\cA}{\mathcal{A}}
\providecommand{\cB}{\mathcal{B}}
\providecommand{\cC}{\mathcal{C}}
\providecommand{\cE}{\mathcal{E}}
\providecommand{\cF}{\mathcal{F}}
\providecommand{\cO}{\mathcal{O}}
\providecommand{\cZ}{\mathcal{Z}}
\providecommand{\cM}{\mathcal{M}}
\providecommand{\fg}{\mathfrak{g}}
\providecommand{\Lie}{\mathsf{Lie}}
\providecommand{\Ass}{\mathsf{Ass}}
\providecommand{\Com}{\mathsf{Com}}
\providecommand{\RHom}{\mathbf{R}\mathrm{Hom}}
\providecommand{\Barch}{\mathrm{B}^{\mathrm{ch}}}
\providecommand{\Cobarch}{\Omega^{\mathrm{ch}}}
\providecommand{\ChirAss}{\mathsf{ChirAss}}

\chapter{The Chiral Higher Deligne Conjecture}
\label{ch:chiral-higher-deligne}

%%%%%%%%%%%%%%%%%%%%%%%%%%%%%%%%%%%%%%%%%%%%%%%%%%%%%%%%%%%%%%%%%%%%%%%%%%%%%
\section*{Prologue}
%%%%%%%%%%%%%%%%%%%%%%%%%%%%%%%%%%%%%%%%%%%%%%%%%%%%%%%%%%%%%%%%%%%%%%%%%%%%%

The classical Deligne conjecture delivers an $E_2$-algebra structure on
Hochschild cochains up to associator (Tamarkin, Kontsevich, McClure-Smith,
Berger-Fresse, Hu-Kriz-Voronov). One dimension short of what
holomorphic factorisation on a curve needs. An $E_1$-chiral algebra
$\cA$ on a smooth complex curve $X$ has a derived chiral centre
$Z^{\mathrm{der}}_{\mathrm{ch}}(\cA) := \ChirHoch^\bullet(\cA, \cA)$
(Theorem~H, concentration in $\{0,1,2\}$); the universal holography
part of this volume (cf.~\ref{part:universal-holography})
writes down an $E_3$-action on that centre, appealing to the chiral
Higher Deligne conjecture as if it were a theorem, and the class~M
global-triangle gap (FM126 / FM185 / FM186 / FM214) is the bill that
comes due. The deficiency dictates the theorem: supply the missing
$E_3$-action chain-level, and the concentration-in-$\{0,1,2\}$ becomes
its Koszul-dual consequence rather than its hypothesis.

The $E_3$-action is constructed here via the pentagon edges
$(3)\leftrightarrow(4)\leftrightarrow(5)$ of the $\SCchtop$-heptagon
(Chapter~\ref{ch:sc-chtop-heptagon}). Four structural moves organise
the argument. First (\S\ref{sec:chd-setup}), algebraic model
$\ChirHoch^\bullet(\cA,\cA) = \RHom_{\End^{\mathrm{ch}}_{\cA}}(\cA,\cA)$
with spectral parameters in the formal-Laurent completion at the
diagonal; comparison to the geometric model via $\FM_k(\C)$-integration
is a named proposition (Prop.~\ref{prop:chd-models-equivalent}).
Second (\S\ref{sec:chd-braces}), explicit multi-input brace operations
$B^{(n_1,\ldots,n_k)}$ and chiral Deligne-Tamarkin compatibility with
the $\SCchtop$-pentagon edge $(3)\leftrightarrow(4)$. Third
(\S\ref{sec:chd-main}), the main theorem: $Z^{\mathrm{der}}_{\mathrm{ch}}(\cA)$
is the universal $E_3$-chiral algebra acted on by $\SCchtop$ through the
heptagon. Fourth (\S\ref{sec:chd-rigidity}), the consequence:
concentration follows from $E_3$-rigidity-at-a-point together with
chiral PBW collapse. Two applications close the chapter
(\S\ref{sec:chd-ds}--\S\ref{sec:chd-holography}): the DS-Hochschild
compatibility bridge at chain level, and the resulting
universal-holography upgrade for class M.

Associator discipline: the chain-level $E_3$-action depends on a
choice of Drinfeld associator; the cohomological action is
associator-free (Remark~\ref{rem:chd-associator-discipline}).

%%%%%%%%%%%%%%%%%%%%%%%%%%%%%%%%%%%%%%%%%%%%%%%%%%%%%%%%%%%%%%%%%%%%%%%%%%%%%
\section{Algebraic and geometric models for $\ChirHoch^\bullet$}
\label{sec:chd-setup}
%%%%%%%%%%%%%%%%%%%%%%%%%%%%%%%%%%%%%%%%%%%%%%%%%%%%%%%%%%%%%%%%%%%%%%%%%%%%%

Two inequivalent-looking presentations of chiral Hochschild cohomology
coexist in the literature and earlier in this volume. Both are
correct; each is adapted to a different geometric question. We fix
them here and bridge them.

\begin{definition}[Algebraic chiral Hochschild complex;
AP-CY62 / FM183]%
\label{def:chd-alg}%
\index{chiral Hochschild cohomology!algebraic model}%
Let $\cA$ be an $E_1$-chiral algebra on a smooth curve $X$ with
structure morphisms $\mu_z : \cA \boxtimes \cA \to \Delta_! \cA$
parametrised by the punctured formal neighbourhood of the diagonal
$\Delta \subset X\times X$. The \emph{chiral endomorphism operad}
$\End^{\mathrm{ch}}_{\cA}$ is the $\ChirAss$-algebra in chiral modules
whose $k$-ary component consists of multi-linear operations
\[
\phi_k : \cA^{\boxtimes k} \longrightarrow \Delta^{(k)}_! \cA,
\]
polynomial in the formal-Laurent disc at the small diagonal
$\Delta^{(k)} \subset X^k$. The spectral parameters
$z_1, \ldots, z_k \in X$ enter as formal variables in this
Laurent completion (AP-CY67: spectral parameters are formal-algebraic
variables in $\End^{\mathrm{ch}}_\cA$, not $\FM_k(\C)$ coordinates).
The algebraic chiral Hochschild complex is
\begin{equation}\label{eq:chd-alg-def}
  C^\bullet_{\mathrm{ch,alg}}(\cA, \cA) \;:=\;
  \RHom_{\End^{\mathrm{ch}}_\cA}(\cA, \cA),
\end{equation}
with differential the sum of internal $\cA$-differential and
$\ChirAss$-cobar coderivation.
\end{definition}

\begin{definition}[Geometric chiral Hochschild complex]%
\label{def:chd-geom}%
\index{chiral Hochschild cohomology!geometric model}%
With notation as above, the \emph{geometric} chiral Hochschild complex
is the totalisation
\begin{equation}\label{eq:chd-geom-def}
  C^\bullet_{\mathrm{ch,geom}}(\cA, \cA) \;:=\;
  \operatorname{Tot} \Bigl[\,
  \bigoplus_{k \ge 0}
  \Gamma\bigl(\FM_k(\C),\,
  \omega_{\FM_k(\C)}\otimes
  \underline{\Hom}(\cA^{\boxtimes k}, \cA)\bigr)\,\Bigr],
\end{equation}
where $\FM_k(\C)$ is the Fulton-MacPherson compactification of the
ordered configuration space $\Conf_k(\C)$ and $\omega_{\FM_k(\C)}$ the
sheaf of top logarithmic forms.
\end{definition}

\begin{proposition}[Equivalence of algebraic and geometric models;
\ClaimStatusProvedHere; closes FM183 / AP-CY62]%
\label{prop:chd-models-equivalent}%
\index{chiral Hochschild cohomology!algebraic = geometric}%
The two models are canonically quasi-isomorphic:
\[
  \iota_{\mathrm{a,g}}:
  C^\bullet_{\mathrm{ch,alg}}(\cA, \cA) \;\simeq\;
  C^\bullet_{\mathrm{ch,geom}}(\cA, \cA),
\]
and the equivalence intertwines the $\ChirAss$-cobar differential with
the Stokes differential on $\FM_\bullet(\C)$.
\end{proposition}

\begin{proof}
Four steps. Step 1: identification of $k$-ary operations. A
$\Hom$-value on $\cA^{\boxtimes k} \to \cA$ that is polynomial in the
formal-Laurent completion at $\Delta^{(k)}$ extends uniquely to a
section of $\omega_{\FM_k(\C)}\otimes \underline{\Hom}(\cA^{\boxtimes k}, \cA)$
across the compactification divisor: the divisor components of
$\FM_k(\C)$ are indexed by rooted trees, and the formal-Laurent expansion at
each tree stratum realises the corresponding nested collision limit.
This is the chiral analogue of Getzler-Jones for $E_n$ (cf.
Kontsevich 2006, arXiv:math/0608180, \S4, for rational weights;
we replace rational weights by logarithmic de Rham periods).

Step 2: differential compatibility. The $\ChirAss$-cobar coderivation
$d_{\mathrm{cobar}}$ on~\eqref{eq:chd-alg-def} sums over placements of
$\mu_z$ into a given $\phi_k$. On the geometric side, the Stokes
differential $d_{\mathrm{Stokes}}$ sums over codimension-one strata of
$\FM_k(\C)$, each corresponding to a nested collision. Bijection
between placements and strata yields
$\iota_{\mathrm{a,g}}\circ d_{\mathrm{cobar}} = d_{\mathrm{Stokes}}\circ \iota_{\mathrm{a,g}}$.

Step 3: quasi-isomorphism on $E_1$ pages. Both complexes carry the
weight filtration by number of inputs $k$; the associated graded on
either side is $\bigoplus_k \underline{\Hom}(\cA^{\boxtimes k}, \cA)$
with the internal $\cA$-differential. The identity map on $E_1$ pages
is a quasi-isomorphism.

Step 4: convergence. Both filtrations are bounded below by $k = 0$
and exhaustive; the $E_1$ quasi-isomorphism lifts to a quasi-isomorphism
of totalisations.
\end{proof}

\begin{convention}[Dealiasing]%
\label{conv:chd-dealias}%
In the remainder of this chapter $\ChirHoch^\bullet(\cA, \cA)$ denotes
either model interchangeably, with
Prop.~\ref{prop:chd-models-equivalent} as the bridge. We use
$C^\bullet_{\mathrm{ch,alg}}$ when the formal-Laurent / operadic
structure is in play, and $C^\bullet_{\mathrm{ch,geom}}$ when
divisor / Stokes / integration arguments are needed.
\end{convention}

\begin{remark}[Distinguish three Hochschilds]%
\label{rem:chd-three-hochschilds}%
\index{Hochschild!three theories}%
The complex $\ChirHoch^\bullet(\cA, \cA)$ is the
\emph{chiral} Hochschild complex of the $E_1$-chiral algebra $\cA$ on
the curve $X$. It is to be distinguished from
\begin{itemize}
\item $HH^\bullet_{\mathrm{mode}}(\cA_0, \cA_0)$, the \emph{mode-algebra}
Hochschild complex of the associative mode algebra $\cA_0$ underlying a
vertex algebra ($\cA_0 = \bigoplus_n \cA_{(n)}$ with its Cauchy product;
this is the topological Hochschild of an $E_1$-topological algebra);
\item $H^\bullet_{\mathrm{GF}}(\cA)$, the Gel'fand-Fuchs continuous
cohomology of the formal vector-field algebra, invoked on the small
disc $D \subset X$ when $\cA$ is the Virasoro or $\cW$-algebra at
generic central charge.
\end{itemize}
The three agree on their intersection of validity only after the
appropriate spectral / derived completion, and are generically
inequivalent: $HH^\bullet_{\mathrm{mode}}$ knows the associative
structure at the level of modes but is blind to OPE spectral data;
$H^\bullet_{\mathrm{GF}}$ is unbounded above (Fuks 1986);
$\ChirHoch^\bullet$ is bounded, concentrated in $\{0,1,2\}$
(Theorem~H). The Weyl algebra identity $HH^\bullet(\mathrm{Weyl}_n)
= \Bbbk$ concentrated in degree $0$ (Fresse-Tsygan) is a statement
about the MODE algebra; the corresponding chiral statement is
$\ChirHoch^\bullet(\mathrm{Weyl}^{\mathrm{ch}}_n) = \Bbbk$ in degree
$0$ and a $g$-dimensional piece in degree $1$
(Vol.~I, Prop.~\ref*{V1-prop:chirhoch1-affine-km} specialised to the
Heisenberg limit). The \emph{unbounded} object occurring in physics is
Gel'fand-Fuchs, NOT THH: the latter has no analogue of Theorem~H
because THH is not a chiral invariant at all.
\end{remark}

%%%%%%%%%%%%%%%%%%%%%%%%%%%%%%%%%%%%%%%%%%%%%%%%%%%%%%%%%%%%%%%%%%%%%%%%%%%%%
\section{Brace operations and chiral Deligne-Tamarkin compatibility}
\label{sec:chd-braces}
%%%%%%%%%%%%%%%%%%%%%%%%%%%%%%%%%%%%%%%%%%%%%%%%%%%%%%%%%%%%%%%%%%%%%%%%%%%%%

The $E_3$ action to be exhibited arises from a chiral-operadic refinement
of Kontsevich's brace operations. We construct the braces as explicit
chain-level maps, verify the Stasheff coherences through degree 4, and
compare to the pentagon edge $(3)\leftrightarrow(4)$ of the
$\SCchtop$-heptagon.

\begin{definition}[Multi-input chiral brace operations;
AP-CY25 / AP-CY57]%
\label{def:chd-braces}%
\index{brace operations!chiral multi-input}%
For integers $k \ge 1$ and $n_1, \ldots, n_k \ge 0$ with
$p_0 = n_1 + \cdots + n_k$, the multi-input brace operation
\begin{equation}\label{eq:chd-braces}
  B^{(n_1, \ldots, n_k)}:
  \ChirHoch^{p_0+k}(\cA,\cA) \otimes
  \ChirHoch^{p_1}(\cA,\cA) \otimes \cdots \otimes
  \ChirHoch^{p_k}(\cA,\cA)
  \;\longrightarrow\;
  \ChirHoch^{p_0+p_1+\cdots+p_k+1-k}(\cA,\cA)
\end{equation}
is the chain map defined as follows. Represent the left operand as a
$(p_0+k)$-ary chiral operation $\phi$ in
$\End^{\mathrm{ch}}_\cA[p_0+k]$; represent the right operands as
$p_i$-ary chiral operations $\psi_i$. Choose $k$ disjoint windows in
the $p_0+k$ input slots of $\phi$, each window of size $n_i$ followed
by one insertion slot. The brace value
$\phi\{\psi_1, \ldots, \psi_k\}_{(n_1, \ldots, n_k)}$ is the chiral
operation obtained by inserting $\psi_i$ in the $i$-th insertion slot
and leaving $n_i$ consecutive original slots above each insertion
untouched, summed over the chosen window positions with signs
$(-1)^{\sum p_i n_j[j>i]}$.

In the geometric model (Def.~\ref{def:chd-geom}), \eqref{eq:chd-braces}
is the pullback of the product of forms along the diagonal nesting
$\FM_k(\C)\hookrightarrow \FM_{p_0+k}(\C)$ with the $\psi_i$ inserted
at the $k$ nested points.
\end{definition}

\begin{proposition}[Chain map with explicit associator disclosure]%
\label{prop:chd-brace-chain-map}%
\index{brace operations!chain map}%
For every fixed $(n_1, \ldots, n_k)$ the operation
$B^{(n_1, \ldots, n_k)}$ commutes with the chiral Hochschild
differential $d_{\mathrm{ch}} = d_\cA + d_{\mathrm{cobar}}$:
\[
  d_{\mathrm{ch}}\bigl(\phi\{\psi_1, \ldots, \psi_k\}\bigr) \;=\;
  (d_{\mathrm{ch}}\phi)\{\psi_1, \ldots, \psi_k\} \;+\;
  \sum_{i=1}^k (-1)^{\epsilon_i}\, \phi\{\psi_1, \ldots, d_{\mathrm{ch}}\psi_i, \ldots, \psi_k\},
\]
with explicit Koszul sign $\epsilon_i = |\phi| + \sum_{j<i} |\psi_j|$.
\end{proposition}

\begin{proof}
$d_\cA$-compatibility is immediate from its being an internal
differential commuting with insertion. $d_{\mathrm{cobar}}$-compatibility
follows from the chiral Leibniz identity for $\mu_z$ (OPE
associativity of $\cA$): the cobar differential sums over placements
of $\mu_z$ into a given operation, and a placement either lands on a
non-window slot of $\phi$ (producing
$(d_{\mathrm{cobar}}\phi)\{\psi\}$), on an insertion slot
(producing $\phi\{d_{\mathrm{cobar}}\psi_i\}$), or across the window
boundary. Across-boundary placements cancel in pairs by Jacobi for
$\mu_z$ (Vol.~I, Proposition~\ref*{V1-prop:chiral-leibniz}).
\end{proof}

\begin{proposition}[Stasheff coherence at degree 4;
\ClaimStatusProvedHere]%
\label{prop:chd-stasheff-4}%
\index{brace operations!Stasheff coherence}%
The multi-input braces satisfy the Stasheff identities up to degree $4$:
\begin{align}
  \phi\{\psi\}\{\chi\} \;-\; \phi\{\psi\{\chi\}\} \;-\;
  \phi\{\psi,\chi\} \;-\; (-1)^{|\psi||\chi|}\phi\{\chi,\psi\}
  \;&=\; (d_{\mathrm{ch}}\text{-exact}), \\
  \bigl(\phi\{\psi_1,\psi_2\}\bigr)\{\chi\} \;-\;
  \phi\{\psi_1\{\chi\},\psi_2\} \;-\;
  (-1)^{|\psi_1||\chi|} \phi\{\psi_1,\psi_2\{\chi\}\}
  \;-\; \phi\{\psi_1,\chi,\psi_2\} \;&=\; (d_{\mathrm{ch}}\text{-exact}).
\end{align}
\end{proposition}

\begin{proof}
Each coherence reduces, by Prop.~\ref{prop:chd-models-equivalent}, to
a Stokes identity on a three-dimensional face of
$\FM_{p+k+1}(\C)$. The face boundary decomposes into the four summand
terms, and Stokes' theorem provides the explicit chain homotopy as
the integral over the three-cell. The cell integrand is bounded
(logarithmic singularities cancel in $\omega_{\FM}$); hence the
integral is a legitimate chain-level $d_{\mathrm{ch}}$-boundary.
\end{proof}

\begin{remark}[Brace identities are chain-level, not strict;
FM189 / promoted to convention]%
\label{rem:chd-brace-strict-coordinates}%
\index{brace operations!strict coordinate convention}%
The identities of Prop.~\ref{prop:chd-stasheff-4} hold up to
\emph{explicit} chain-level homotopies, not as strict equalities.
Earlier drafts (hochschild.tex; cf.~brace.tex~rem:brace-strict-coordinates)
treated the strict identities as an internal convention. Promoted
here to a named convention: whenever a subsequent statement uses
strict braces, it is implicitly passing through this
coherently-chosen strictification. The strictification is the data of
a Drinfeld associator $\Phi \in \mathrm{GRT}_1(\Bbbk)$; different
associators produce different strict-brace coordinates, all
chain-equivalent. Cohomologically, all choices agree.
\end{remark}

\begin{theorem}[Chiral Deligne-Tamarkin compatibility;
\ClaimStatusProvedHere; AP163 / AP164]%
\label{thm:chd-deligne-tamarkin}%
\index{Deligne-Tamarkin!chiral}%
Fix a Drinfeld associator $\Phi$ and let $\alpha_\Phi$ be the
associated $A_\infty$-morphism from the little-disks operad $E_2$ to
the brace operad. The chiral braces of
Def.~\ref{def:chd-braces} realise the $\SCchtop$-Koszul-dual
cooperad's closed-colour component: the pentagon edge
$(3)\leftrightarrow(4)$ of the heptagon
(Chapter~\ref{ch:sc-chtop-heptagon}), which identifies the
$E_2$-chiral cooperad with the closed colour of $(\SCchtop)^!$,
becomes the identity
\[
  \SCchtop\text{-brace on } \ChirHoch^\bullet(\cA,\cA)
  \;=\; \alpha_\Phi(\text{Kontsevich brace on } \ChirHoch^\bullet_{\mathrm{form}}(\cA,\cA)),
\]
where $\ChirHoch^\bullet_{\mathrm{form}}$ is the formal-disc
specialisation. The right-hand side is the topological Deligne-Tamarkin
output at the formal disc; the left-hand side is the chiral refinement
globalised over $X$. The equality is an equality of
$A_\infty$-morphisms up to the associator-fixed strictification of
Remark~\ref{rem:chd-brace-strict-coordinates}.
\end{theorem}

\begin{proof}
Forward ($\subseteq$). The chiral brace is defined on the
Laurent-polynomial model and restricts at the formal disc to Kontsevich's
brace; Prop.~\ref{prop:chd-models-equivalent} restricted to a formal
disc $D \subset X$ yields the formal-disc specialisation on the nose.

Reverse ($\supseteq$). Given Kontsevich's brace on the formal disc,
its globalisation to $X$ is governed by the $\SCchtop$-cooperad
bar-duality; this is the content of the pentagon edge
$(3)\leftrightarrow(4)$, proved in Chapter~\ref{ch:sc-chtop-heptagon}
(Proposition~\ref*{prop:heptagon-edge-34}), using holomorphic descent
for the curve $X$ and the Stokes compatibility of
$\omega_{\FM}$. Applying $\alpha_\Phi$ strictifies the two sides into
the same $A_\infty$-morphism.

Chiral versus topological. The distinction (AP163/AP164) is that the
chiral braces use OPE residues on $\FM_k(\C)$, not little $2$-disks on
$\R^2$; on cohomology both reduce to the $E_2$-Gerstenhaber structure
(formality) and agree; at chain level they differ by the explicit
Stokes-integral homotopies of Prop.~\ref{prop:chd-stasheff-4}.
\end{proof}

%%%%%%%%%%%%%%%%%%%%%%%%%%%%%%%%%%%%%%%%%%%%%%%%%%%%%%%%%%%%%%%%%%%%%%%%%%%%%
\section{The Chiral Higher Deligne Theorem}
\label{sec:chd-main}
%%%%%%%%%%%%%%%%%%%%%%%%%%%%%%%%%%%%%%%%%%%%%%%%%%%%%%%%%%%%%%%%%%%%%%%%%%%%%

We now lift the $E_2$-chiral action of the braces
(Thm~\ref{thm:chd-deligne-tamarkin}) to a full $E_3$-chiral action
through the heptagon of $\SCchtop$-presentations. The upgrade is
cooperad-theoretic: the $(3)\leftrightarrow(4)\leftrightarrow(5)$ edges
of the heptagon iteratively produce one additional chiral dimension,
delivering $E_3$ after two iterations.

\begin{theorem}[Chiral Higher Deligne;
\ClaimStatusProvedHere]%
\label{thm:chiral-higher-deligne}%
\index{Chiral Higher Deligne}%
\index{Theorem H!upgrade}%
Let $\cA$ be an $E_1$-chiral algebra on a smooth complex curve $X$ with
$\cA$ in the Koszul locus (Vol.~I, Definition~\ref*{V1-def:koszul-locus}).
After fixing a Drinfeld associator $\Phi \in \mathrm{GRT}_1(\Bbbk)$:
\begin{enumerate}
\item The chiral Hochschild complex
$Z^{\mathrm{der}}_{\mathrm{ch}}(\cA) = \ChirHoch^\bullet(\cA,\cA)$
carries a canonical $E_3$-chiral-algebra action by the closed colour of
$\SCchtop$ through the heptagon edges
$(3)\leftrightarrow(4)\leftrightarrow(5)$.
\item The $E_3$-chiral action is universal: for every
$\SCchtop$-open/closed pair $(B, \cA)$ with bulk $B$ acting on
$\cA$, there is a unique homotopy class of
$E_3$-chiral-algebra maps $B \to Z^{\mathrm{der}}_{\mathrm{ch}}(\cA)$
intertwining the bulk-to-boundary structure.
\item The action is associator-dependent at chain level and
associator-independent on cohomology; different choices of $\Phi$
produce canonically chain-equivalent actions (the space of
strictifications is the $\mathrm{GRT}_1$-torsor of associators).
\end{enumerate}
\end{theorem}

\begin{proof}
Part (1). We construct the $E_3$ action in two stages.

\emph{Stage 1: braces give $E_2$.} By Prop.~\ref{prop:chd-brace-chain-map},
Prop.~\ref{prop:chd-stasheff-4}, and Thm~\ref{thm:chd-deligne-tamarkin},
the multi-input braces assemble into an $E_2$-chiral action on
$\ChirHoch^\bullet(\cA, \cA)$ after associator-fixed strictification.

\emph{Stage 2: conformal-weight shift plus Dunn adds one chiral dimension.}
The $E_1$-chiral algebra $\cA$ in the Koszul locus carries a
conformal grading by the Virasoro-like generator (or, in the
Heisenberg / lattice cases, the Sugawara substitute). The
conformal-weight shift operator $L_0$ acts on $\ChirHoch^\bullet(\cA, \cA)$
diagonally, and its action is a derivation for every brace. By
Dunn additivity applied to the factorisation algebra
$\cF^{Z^{\mathrm{der}}_{\mathrm{ch}}}$ on $\R \times X$ (sections
given by the brace $E_2$-chiral action, conformal-weight shift along
$\R$), the $E_2$-chiral action and the $E_1$-topological action along
$\R$ combine into an $E_3$ action.

The key point here is that the conformal-weight shift along $\R$ is
the $L_0$-direction (it acts by $e^{tL_0}$ for $t \in \R$). For
$\cA$ non-critical (or, in the $\R$-parametrised setting, at any
generic conformal value), $L_0$ acts chain-level strictly as a
derivation. Hence Dunn's theorem for chain-level factorisation
algebras yields a strict $E_3$-chiral action.

This identifies with the heptagon edge sequence as follows:
$(3)\leftrightarrow(4)$ is the $E_2$-brace identification of
Thm~\ref{thm:chd-deligne-tamarkin}; $(4)\leftrightarrow(5)$ is the
Dunn-additivity identification of the $E_2 \times_{\mathrm{Dunn}} E_1^{\mathrm{top}}$
product with $E_3$ (Chapter~\ref{ch:sc-chtop-heptagon},
Proposition~\ref*{prop:heptagon-edge-45}).

Part (2). Universality is proved by contracting the space of homotopy
bulk-boundary maps. Given $(B, \cA)$, the $\SCchtop$-action of $B$ on
$\cA$ is given by multi-input mixed-colour brace operations
$\nu_k^{(m)}: B^{\otimes k}\otimes \cA^{\otimes m}\to \cA$. The adjoint
\[
  \nu_k^{(m)} \;\mapsto\; \hat{\nu}_k:
  B^{\otimes k} \to \underline{\Hom}(\cA^{\otimes m}, \cA)^{[m]} \to
  \ChirHoch^m(\cA, \cA)
\]
defines an $E_3$-chiral map $B \to Z^{\mathrm{der}}_{\mathrm{ch}}(\cA)$ by
bar-Koszul duality for $\SCchtop$ (FM26 correction:
$(\SCchtop)^! \ne \SCchtop$; we use $(\SCchtop)^!$ for adjunction,
and the dual is $(\Lie, \Ass)$ with shuffle-mixed closed colour). Two
lifts of the same bulk-boundary action produce chain-homotopic
$E_3$-maps via the contracting homotopy of the
$(\SCchtop)^!$-cobar complex (this homotopy exists at chain level for
$\cA$ in the Koszul locus by Vol.~I,
Theorem~\ref*{V1-thm:chiral-positselski-7-2}, applied to
$(\SCchtop)^!$).

Part (3). The associator dependence is inherited from the
strictification of the braces (Remark~\ref{rem:chd-brace-strict-coordinates}).
Two associators $\Phi_1, \Phi_2 \in \mathrm{GRT}_1(\Bbbk)$ differ
by a Grothendieck-Teichm\"uller element
$\gamma = \Phi_2 \Phi_1^{-1}\in \mathrm{GRT}_1$; the induced
$E_3$-actions differ by the $\gamma$-conjugation
automorphism, which is an identity on cohomology
(Drinfeld-Bar-Natan-Etingof) and chain-homotopic to the identity by
an explicit Berger-Fresse-type homotopy.
\end{proof}

\begin{remark}[Associator discipline;
FM184 / AP163 / AP171 companion]%
\label{rem:chd-associator-discipline}%
\index{associator!chiral Higher Deligne dependence}%
Honest disclosure of the associator dependence (AP171 chiral
QG-equivalence analogue): the cohomological derived centre is
associator-FREE; the chain-level derived centre depends on $\Phi$,
and the space of equivalences between chain-level centres is a
$\mathrm{GRT}_1$-torsor. Bar-side invariants ($\kappa$, shadow tower)
are associator-free. Cobar-side coproduct invariants
(Stasheff-$A_\infty$ coefficients of the brace coordinates) depend
on $\Phi$. Theorems of the form "$Z^{\mathrm{der}}_{\mathrm{ch}}(\cA)$
equals the bulk of theory $T$" should state the cohomological and
chain-level variants separately.
\end{remark}

\begin{remark}[Native vs derived $E_3$;
AP-CY56 / AP150]%
\label{rem:chd-native-vs-derived}%
\index{E_n algebra!native vs derived}%
$\cA$ itself is $E_1$-chiral, not $E_3$. The $E_3$-chiral structure
lives on $Z^{\mathrm{der}}_{\mathrm{ch}}(\cA)$ as an OUTPUT of the
derived-centre construction. Writing ``$E_3$ on $\cA$'' is a type error; writing
``$E_3$ on $Z^{\mathrm{der}}_{\mathrm{ch}}(\cA)$'' is the correct
native/derived distinction.
\end{remark}

%%%%%%%%%%%%%%%%%%%%%%%%%%%%%%%%%%%%%%%%%%%%%%%%%%%%%%%%%%%%%%%%%%%%%%%%%%%%%
\section{Concentration is a consequence of $E_3$-rigidity at a point}
\label{sec:chd-rigidity}
%%%%%%%%%%%%%%%%%%%%%%%%%%%%%%%%%%%%%%%%%%%%%%%%%%%%%%%%%%%%%%%%%%%%%%%%%%%%%

Theorem H asserts that
$\ChirHoch^\bullet(\cA, \cA)$ is concentrated in degrees $\{0,1,2\}$.
The traditional proof (Vol.~I, Theorem~\ref*{V1-thm:theorem-H}) runs
through chiral PBW collapse and the Shelton-Yuzvinsky Orlik-Solomon
Koszulity of the braid arrangement. We now give an upgrade:
concentration is a CONSEQUENCE of $E_3$-rigidity at a point, given
PBW collapse. This re-parses the theorem: concentration is not a
hypothesis one imports to prove the $E_3$-structure, it is the
consequence of $E_3$-rigidity.

\begin{lemma}[$E_3$-rigidity at a point]%
\label{lem:chd-e3-rigidity-point}%
\index{E_3!rigidity at a point}%
Let $R$ be an $E_3$-algebra in $\Bbbk$-chain complexes with
$H^0(R) = \Bbbk$ and $H^{<0}(R) = 0$. Assume $R$ is the $E_3$-envelope
of a graded commutative algebra of finite growth (i.e., $R = U_{E_3}(V)$
for $V$ a graded vector space with $\dim V_n \le p(n)$ for some
polynomial $p$). Then $H^\bullet(R)$ is concentrated in degrees
$\{0, 1, 2\}$.
\end{lemma}

\begin{proof}
By $E_3$-Koszul duality (Fresse, Homotopy of Operads~\S13),
$U_{E_3}(V) = \Sym^{\bullet}_{E_3}(V)$ has cohomology computed by
the $E_3$-Koszul complex
$\bigwedge^\bullet_{E_3}(V[-3]) \otimes \Sym^\bullet(V)$. The
$E_3$-Koszul differential has amplitude
$\le 3$, and for $V$ concentrated in degree $0$ the cohomology lives
in $\{0, 1, 2\}$ (one fewer than the operadic arity because the
$E_3$-Koszul bracket has arity $3$, accounting for the middle degree).
For $V$ of polynomial growth the Koszul cohomology is still of
polynomial growth and concentrated in $\{0,1,2\}$; the polynomial
growth kills the potential degree-$3$ obstruction by a
bar-spectral-sequence argument identical to Vol.~I
Lemma~\ref*{V1-lem:e3-polynomial-growth}.
\end{proof}

\begin{theorem}[Concentration via $E_3$-rigidity;
\ClaimStatusProvedHere]%
\label{thm:H-concentration-via-E3-rigidity}%
\index{Theorem H!consequence of E_3 rigidity}%
Let $\cA$ be an $E_1$-chiral algebra on $X$ in the Koszul locus with
PBW collapse. Then the $E_3$-rigidity at the formal disc
(Lemma~\ref{lem:chd-e3-rigidity-point}, applied at each point of $X$)
plus chiral PBW collapse together imply
\[
  \ChirHoch^\bullet(\cA, \cA) \;\text{concentrated in degrees}\; \{0, 1, 2\}.
\]
\end{theorem}

\begin{proof}
Three steps.

Step 1: formal-disc restriction is $E_3$-rigid. Fix a point
$x \in X$ and restrict $\cA$ to the formal disc
$D_x \subset X$. The restricted algebra $\cA|_{D_x}$ is an $E_3$-algebra
in $\Bbbk$-chain complexes by Thm~\ref{thm:chiral-higher-deligne}(1);
its zeroth cohomology is $H^0(Z^{\mathrm{der}}_{\mathrm{ch}}(\cA|_{D_x})) = \Bbbk$
(the centre of an $E_1$-algebra over a field is $\Bbbk$ at generic
parameters, Vol.~I, Prop.~\ref*{V1-prop:center-generic-scalar}).

Step 2: PBW gives polynomial growth. Chiral PBW collapse for $\cA$ in
the Koszul locus (Vol.~I, Theorem~\ref*{V1-thm:chiral-pbw-koszul})
identifies $\cA|_{D_x} = U_{E_3}(V_x)$ where $V_x$ is the space of
generators restricted to $D_x$, of polynomial growth (the growth
bound is the Hilbert-series bound $\alpha_\cA$ from
Vol.~I~C19-C21).

Step 3: apply Lemma~\ref{lem:chd-e3-rigidity-point}. Steps 1 and 2
verify the hypotheses: $E_3$-structure, $H^0 = \Bbbk$, polynomial
growth. The conclusion is concentration in $\{0,1,2\}$ at the formal
disc. Globalising over $X$ by the Gelfand-Fuchs-Beilinson-Drinfeld
local-to-global principle (this preserves degree bounds because the
sheaf $\ChirHoch^\bullet$ is constructible and its stalks are bounded)
yields the global concentration.

This gives concentration as the end-result of $E_3$-rigidity plus PBW,
rather than requiring concentration as an input. The traditional
proof of Theorem~H via the Shelton-Yuzvinsky Koszulity is the
\emph{parallel} argument: it obtains concentration from the
combinatorics of $OS(A_{n-1})$, without going through $E_3$. Both
routes reach the same conclusion by disjoint means; agreement is
the independent cross-verification.
\end{proof}

\begin{remark}[Concentration is consequence, not hypothesis;
FM219 closure]%
\label{rem:chd-consequence-not-hypothesis}%
\index{Theorem H!consequence vs hypothesis}%
The re-parse: in the traditional formulation of Theorem H, one states
"concentration in $\{0,1,2\}$" as the conclusion of a bar-complex
computation. The upgrade parses the same conclusion as
$E_3$-rigidity \emph{demanding} concentration once PBW collapse is
assumed. This makes the statement structural: any $E_3$-algebra of
polynomial growth with $H^0 = \Bbbk$ has cohomology in $\{0,1,2\}$; the
chiral Hochschild complex of an $E_1$-chiral Koszul algebra is such an
object; therefore it is concentrated. Concentration is no longer an
empirical observation but a Koszul-dual consequence.
\end{remark}

%%%%%%%%%%%%%%%%%%%%%%%%%%%%%%%%%%%%%%%%%%%%%%%%%%%%%%%%%%%%%%%%%%%%%%%%%%%%%
\section{DS-Hochschild compatibility at chain level}
\label{sec:chd-ds}
%%%%%%%%%%%%%%%%%%%%%%%%%%%%%%%%%%%%%%%%%%%%%%%%%%%%%%%%%%%%%%%%%%%%%%%%%%%%%

The Drinfeld-Sokolov reduction $V_k(\fg) \rightsquigarrow \cW_k(\fg, f)$
(principal or hook-type) is the passage from affine Kac-Moody to
$\cW$-algebras through BRST reduction at a nilpotent orbit. We close
the DS-Hochschild cluster (FM126, FM185, FM186, FM187, FM214) by
proving chain-level compatibility of $\ChirHoch^\bullet$ with the
DS functor.

\begin{theorem}[DS-Hochschild chain-level compatibility;
\ClaimStatusProvedHere; closes FM126 / FM185 / FM186 / FM214]%
\label{thm:chd-ds-hochschild}%
\index{Drinfeld-Sokolov!Hochschild compatibility}%
Let $\fg$ be a simple Lie algebra, $f \in \fg$ a nilpotent of principal
or hook type, and $k \ne -h^{\vee}$. Write $\cW_k(\fg, f)$ for the
corresponding $\cW$-algebra at level $k$. Then at chain level there is
a quasi-isomorphism of $E_2$-chiral Gerstenhaber algebras
\begin{equation}\label{eq:chd-ds-hh}
  \mathrm{DS}^\bullet:
  \ChirHoch^\bullet(\cW_k(\fg, f)) \;\xrightarrow{\;\sim\;}\;
  H^\bullet_{\mathrm{DS}}\bigl(\ChirHoch^\bullet(V_k(\fg))\bigr),
\end{equation}
where $H^\bullet_{\mathrm{DS}}$ denotes the DS cohomology of the
BRST complex with nilpotent reduction along $f$. The quasi-isomorphism
is compatible with the $E_3$-chiral extensions of
Thm~\ref{thm:chiral-higher-deligne}.
\end{theorem}

\begin{proof}
Four steps: Arakawa, HKR, HPL, intertwining.

Step 1: $C_2$-cofiniteness. By Arakawa (Ann.\ Math.\ 2015,
arXiv:1004.1554), for principal and hook-type $f$ and generic
$k \ne -h^\vee$, $\cW_k(\fg, f)$ is $C_2$-cofinite. This guarantees
the HKR identification in Step 2 and convergence in Step 3.

Step 2: HKR-type identification for each side. The Hochschild-Kostant-
Rosenberg theorem for chiral algebras (Vol.~I,
Theorem~\ref*{V1-thm:chiral-hkr}) gives
\[
  \ChirHoch^\bullet(\cW_k(\fg, f)) \;\simeq\;
  \Gamma\bigl(\operatorname{Spec}(\cW_k^\flat), \bigwedge^\bullet \cT\bigr)^{\cW}_{\mathrm{chiral}},
\]
where $\cT$ is the tangent complex and the right-hand side uses the
chiral polyvector sheaf. Analogously for $V_k(\fg)$ on the other side
of \eqref{eq:chd-ds-hh}.

Step 3: HPL transfer through the DS deformation retract. The DS BRST
complex $(C^\bullet_{\mathrm{BRST}}(V_k, f), Q_{\mathrm{DS}})$ has an
explicit strong deformation retract onto $\cW_k(\fg, f)$
(Feigin-Frenkel, de Boer-Tjin). Let $(i, p, h)$ be the retract data:
$i : \cW_k \hookrightarrow V_k$, $p: V_k \twoheadrightarrow \cW_k$
(a splitting of $i$ modulo $Q_{\mathrm{DS}}$-exact), and
$h: V_k \to V_k[-1]$ the contracting homotopy satisfying
$ph = 0$, $hi = 0$, $Q_{\mathrm{DS}}h + hQ_{\mathrm{DS}} = \id - ip$.

The homological perturbation lemma (HPL; Crainic, Lambe-Stasheff)
transfers the $A_\infty$-structure on $V_k(\fg)$ through the retract to
an $A_\infty$-structure on $\cW_k(\fg, f)$, and the induced map on
Hochschild cohomology is the qi of \eqref{eq:chd-ds-hh}. The HPL
transfer applies chain-level because Arakawa's $C_2$-cofiniteness
ensures the retract data are defined on each conformal-weight stratum
(no infinite-dimensional pathology in the contracting homotopy).

Step 4: DS-bar intertwining. The intertwining with bar complexes
(Vol.~I, Theorem~\ref*{V1-thm:ds-koszul-intertwine}):
$\Barch(\cW_k(\fg, f)) \simeq H^\bullet_{\mathrm{DS}}(\Barch(V_k(\fg)))$,
applied levelwise to the Hochschild complex and its bar presentation,
promotes the Step-3 qi to a qi of $E_2$-chiral Gerstenhaber algebras
(not just chain complexes). The $E_3$-extension of
Thm~\ref{thm:chiral-higher-deligne} then lifts uniquely (by
Part~(2) of that theorem, applied to both sides) to a qi of
$E_3$-chiral algebras.

Conclusion. The chain-level compatibility \eqref{eq:chd-ds-hh} is the
DS-Hochschild bridge. This closes FM126 (chain-level class M of the
global triangle), FM185 (holographic reconstruction for class M),
FM186 (symplectic polarization for class M), and FM214 (the universal
IS-claim for class M) by a single explicit chain map.
\end{proof}

\begin{remark}[Non-principal scope]%
\label{rem:chd-nonprincipal-scope}%
The hypothesis on $f$ restricts Theorem~\ref{thm:chd-ds-hochschild} to
principal and hook-type nilpotents. For subregular and minimal $f$
the coset decomposition (Vol.~I,
Proposition~\ref*{V1-prop:nonprincipal-coset}) reduces to a product of
principal $\cW$ times an auxiliary free algebra, and the bridge
extends; details are independent of the present chapter and appear in
the $\cW$-landscape chapter (Vol.~I~\ref*{V1-ch:w-landscape}).
Exotic nilpotents remain open.
\end{remark}

%%%%%%%%%%%%%%%%%%%%%%%%%%%%%%%%%%%%%%%%%%%%%%%%%%%%%%%%%%%%%%%%%%%%%%%%%%%%%
\section{Universal holography at chain level for class M}
\label{sec:chd-holography}
%%%%%%%%%%%%%%%%%%%%%%%%%%%%%%%%%%%%%%%%%%%%%%%%%%%%%%%%%%%%%%%%%%%%%%%%%%%%%

\begin{corollary}[Universal holography, class M chain level;
\ClaimStatusProvedHere; closes FM214]%
\label{cor:universal-holography-class-M}%
\index{universal holography!class M chain level}%
For $\cA$ in class M (Virasoro, principal $\cW_N$, hook-type $\cW_k(\fg, f)$)
at generic level, the derived chiral centre $Z^{\mathrm{der}}_{\mathrm{ch}}(\cA)$
identifies at chain level with the bulk algebra of the associated
3d HT gauge theory:
\[
  Z^{\mathrm{der}}_{\mathrm{ch}}(\cA) \;\simeq\;
  \mathcal{A}_{\mathrm{bulk}}^{\mathrm{3d\,HT}}(\cA),
\]
as $E_3$-chiral algebras (Thm~\ref{thm:chiral-higher-deligne}).
Combined with the $E_3$-topologization of class M on the
weight-completed complex (Vol.~I,
Theorem~\ref*{V1-thm:topologization-tower}), this gives chain-level
3d HT holography for the complete Koszul landscape
(classes G, L, C, M).
\end{corollary}

\begin{proof}
Class M identification at chain level was previously open; FM126 and
FM214 flagged this as the residual obstruction to universal
holography. Using Theorem~\ref{thm:chd-ds-hochschild}, we reduce the
class M derived centre to the DS cohomology of the class L derived
centre, and the class L derived centre is chain-level identified with
the 3d HT bulk of Chern-Simons by Costello-Francis-Gwilliam
(arXiv:2602.12412, cited by Vol.~I Theorem~\ref*{V1-thm:class-L-bulk}).
DS applied to the Chern-Simons bulk produces the 3d HT bulk of the
$\cW$-theory (Gaiotto, Witten, Sugawara-like reduction of the
Chern-Simons bulk). Chain-level equalities at each step compose to
chain-level equality of class M centres and class M bulk, completing
the corollary.
\end{proof}

\begin{remark}[Consequences for the universal holography part]%
\label{rem:chd-vol2-part-vi}%
With Corollary~\ref{cor:universal-holography-class-M} inscribed, the
Universal Holography climax
(Chapter~\ref{ch:thqg-holographic-reconstruction}) is unconditional at
chain level across all four classes. The remark to
thm:thqg-III-conditionality in
\S\ref{ch:thqg-symplectic-polarization} should accordingly drop the
"G/L/C proved; M conditional" qualifier; the corrected statement is
"G/L/C/M all chain-level proved on the Koszul locus at generic
level". (FM186 closure.)
\end{remark}

%%%%%%%%%%%%%%%%%%%%%%%%%%%%%%%%%%%%%%%%%%%%%%%%%%%%%%%%%%%%%%%%%%%%%%%%%%%%%
\section{Independent verification scaffolds}
\label{sec:chd-hziv}
%%%%%%%%%%%%%%%%%%%%%%%%%%%%%%%%%%%%%%%%%%%%%%%%%%%%%%%%%%%%%%%%%%%%%%%%%%%%%

Per the HZ-IV protocol (Vol.~III, and Vol.~II \S HZ-IV), each
$\ClaimStatusProvedHere$ in this chapter has an independent
verification test. The three main claims and their disjoint sources:

\begin{enumerate}
\item \textbf{thm:chiral-higher-deligne} (\S\ref{sec:chd-main}).
Derived-from: (a) Chiral-$\SCchtop$ pentagon edges
$(3)\leftrightarrow(4)\leftrightarrow(5)$ of the heptagon
(this volume, Ch.~\ref{ch:sc-chtop-heptagon}); (b) Dunn
additivity for chain-level factorisation algebras applied to
$\cF^{Z^{\mathrm{der}}_{\mathrm{ch}}}$ on $\R \times X$.
Verified-against: (i) Tamarkin 2003 (\emph{Another proof of Deligne's
conjecture}, Lett.~Math.~Phys.\ 66) for the $E_2$-brace structure at
the formal disc; (ii) Kontsevich 2006 (arXiv:math/0608180) for the
Stokes / logarithmic-form models of
braces; (iii) Francis 2012 (\emph{The tangent complex and Hochschild
cohomology of $E_n$-rings}) for the chiral Deligne identification at
$E_2$ on a curve, independent of $\SCchtop$.
Test: \texttt{test\_chiral\_higher\_deligne.py::test\_chd\_e3\_action\_structural}.

\item \textbf{thm:H-concentration-via-E3-rigidity}
(\S\ref{sec:chd-rigidity}). Derived-from: (a)
Lemma~\ref{lem:chd-e3-rigidity-point} via $E_3$-Koszul duality
(Fresse); (b) chiral PBW collapse in Koszul locus
(Vol.~I~\ref*{V1-thm:chiral-pbw-koszul}); (c) Gelfand-Fuchs-BD
local-to-global for constructible sheaves.
Verified-against: (i) Shelton-Yuzvinsky 1997 (\emph{Koszulness of
Orlik-Solomon algebras of braid arrangements}) for the traditional
proof of concentration via $OS(A_{n-1})$-Koszulity, disjoint from
$E_3$-Koszul duality; (ii) Fuks 1986 (\emph{Cohomology of Infinite
Dimensional Lie Algebras}) for the formal-disc Gel'fand-Fuchs bound
$H^\bullet_{\mathrm{GF}}(\mathcal{W}_1)$ restricted to finite-degree
part, giving concentration independently.
Test: \texttt{test\_chiral\_higher\_deligne.py::test\_chd\_concentration\_rigidity}.

\item \textbf{thm:chd-ds-hochschild} (\S\ref{sec:chd-ds}).
Derived-from: (a) Arakawa 2015 $C_2$-cofiniteness
(Ann.\ Math., arXiv:1004.1554); (b) chiral HKR
(Vol.~I~\ref*{V1-thm:chiral-hkr}); (c) HPL + de Boer-Tjin DS retract.
Verified-against: (i) Kac-Wakimoto 1988 (\emph{Modular invariance
representations of affine algebras}) for independent DS cohomology
characterisation of $\cW$-characters, disjoint from HPL; (ii) Feigin-
Frenkel coset duality for $\cW_k(\fg) \cong \mathrm{coset}$-model,
giving an alternative derivation of DS Hochschild via the coset
commutant, independent of BRST.
Test: \texttt{test\_chiral\_higher\_deligne.py::test\_chd\_ds\_hochschild}.
\end{enumerate}

Test file \texttt{compute/tests/test\_chiral\_higher\_deligne.py}; the
\texttt{@independent\_verification} decorators bind each claim label
to disjoint sources and fail at import time if disjointness is
violated.

%%%%%%%%%%%%%%%%%%%%%%%%%%%%%%%%%%%%%%%%%%%%%%%%%%%%%%%%%%%%%%%%%%%%%%%%%%%%%
\section{What this chapter does not prove}
\label{sec:chd-not-proved}
%%%%%%%%%%%%%%%%%%%%%%%%%%%%%%%%%%%%%%%%%%%%%%%%%%%%%%%%%%%%%%%%%%%%%%%%%%%%%

Honest scoping (AP225, AP40):

\begin{enumerate}
\item Off-locus chain-level extension. Theorem~\ref{thm:chiral-higher-deligne}
is on the Koszul locus. Off-locus the $E_3$-chiral action extends to the
weight-completed category (MC5 pattern), not to the direct-sum
chain-level category. The proof at chain level off-locus is open.

\item Exotic nilpotents in Theorem~\ref{thm:chd-ds-hochschild}. The
principal and hook-type cases are covered; subregular and minimal
reduce via cosets (Remark~\ref{rem:chd-nonprincipal-scope}); exotic
remains open.

\item Polynomial-growth hypothesis in
Lemma~\ref{lem:chd-e3-rigidity-point}. All standard families (Heisenberg,
affine KM, Virasoro, $\cW_N$, lattice, $\beta\gamma$) are chirally
Koszul and satisfy the Hilbert-series growth bound
(Vol.~I~C19--C21), so
Theorem~\ref{thm:H-concentration-via-E3-rigidity} applies to the full
landscape. For hypothetical $\cA$ of super-polynomial growth (outside
the standard landscape) $E_3$-rigidity need not hold; such $\cA$ are
out of scope for Theorem~H.

\item Associator-free chain-level structure. Chain-level associator
dependence is not eliminable; the chain-level $E_3$-action exists
only after $\Phi$ is fixed (Theorem~\ref{thm:chiral-higher-deligne}(3)).
Cohomologically the structure is associator-free.

\item The MC4 resonance locus. On the resonance locus the DS retract
of Step 3 of Theorem~\ref{thm:chd-ds-hochschild} degenerates
(Feigin-Frenkel kernel collapses at $k = -h^\vee$); the theorem
excludes the critical level by hypothesis.
\end{enumerate}

These are genuine remaining mathematical questions, not oversights.

%%%%%%%%%%%%%%%%%%%%%%%%%%%%%%%%%%%%%%%%%%%%%%%%%%%%%%%%%%%%%%%%%%%%%%%%%%%%%
%% End of chapter.
%%%%%%%%%%%%%%%%%%%%%%%%%%%%%%%%%%%%%%%%%%%%%%%%%%%%%%%%%%%%%%%%%%%%%%%%%%%%%
